% Tabla: posición de los venues (cuartil SCImago para revistas, rango ICORE para conferencias).
\begin{table}[H]
\caption{Posición de los venues de los nueve artículos}\label{tab:ranking}
\small
\begin{tabularx}{\textwidth}{@{}L{3.6cm}YL{3.4cm}@{}}
\toprule
\textbf{Artículo} & \textbf{Venue} & \textbf{Indicador} \\
\midrule
\textcite{nigro2022}      & \textit{Algorithms} (MDPI), revista                       & SJR Q2 \\
\textcite{kwedlo2021}     & \textit{IEEE Access} (IEEE), revista                      & SJR Q1 \\
\textcite{kwedlo2024}     & ICPP (ACM), conferencia                                   & ICORE B \\
\addlinespace
\textcite{hess2023}       & AISTATS (PMLR), conferencia                               & ICORE A \\
\textcite{bellavita2026}  & IPDPS (IEEE), conferencia                                 & ICORE A \\
\textcite{xu2021}         & \textit{Computer Networks} (Elsevier), revista            & SJR Q1 \\
\addlinespace
\textcite{li2023}         & \textit{Data Mining and Knowledge Discovery} (Springer), revista & SJR Q1 \\
\textcite{bellavita2025}  & PPoPP (ACM), conferencia                                  & ICORE B \\
\textcite{voevodski2024}  & \textit{ACM Transactions on Knowledge Discovery from Data}, revista & SJR Q1 \\
\bottomrule
\end{tabularx}

{\footnotesize\textit{Nota.} Revistas: mejor cuartil del SCImago Journal Rank 2024, calculado con
datos de Scopus. Conferencias: rango de ICORE 2026 (A* es el más alto, seguido de A, B y C). PPoPP
figuraba como A en CORE 2021 y pasó a B en CORE 2023 porque no presentó datos para su revisión.
Consultado el 13 de septiembre de 2026. Los grupos de filas corresponden a los Temas 1, 2 y 3.\par}
\end{table}
